% Generated by generate.py; do not edit.
\begin{table}
\caption{Every ATL rule, template section and runtime step of the EMF chain that implements the formulation or QAOA, with the reference its comment cites, read from the code.}\label{tab:traceability}
\small
\begin{tabularx}{\linewidth}{@{}l>{\raggedright\arraybackslash}X@{}}
\toprule
Unit & What it implements, as its comment says \\
\midrule
\multicolumn{2}{@{}l}{\textbf{ATL}, \texttt{problem2conflictgraph.atl}: Problem $\to$ Conflict Graph} \\
\texttt{rule Problem2ConflictGraph} & \cite[§5.2, Fig. 3]{Paper}: the conflict graph G, exams as nodes and conflicts as edges. \\
\texttt{rule Problem2ConflictGraph} & \cite[§5.3, Eq. (4)]{Paper}: X has only M columns, so the M batches bound the schedule. \\
\texttt{rule Exam2Exam} & \cite[§5.2, Fig. 3]{Paper}: each exam is a node of G. \\
\texttt{rule Exam2Exam} & \cite[§4.2]{Glover}: the weight w\_j of the set packing objective, the exam's priority (ADR-0009). \\
\texttt{helper sharesStudentWith} & \cite[§5.2, Algorithm 1, Eq. (3)]{Paper}: C[i,j] = 1 iff some student sits both exam i and exam j. \\
\texttt{lazy rule Conflict} & \cite[§5.2, Fig. 3]{Paper}: one edge of G per pair with C[i,j] = 1, emitted once, in problem order. \\
\midrule
\multicolumn{2}{@{}l}{\textbf{Acceleo}, \texttt{main.mtl}: Conflict Graph and Solver Configuration $\to$ \texttt{solve\_<name>.py}} \\
\texttt{CONFLICTS} & \cite[§5.2, Algorithm 1, Eq. (3)]{Paper} \cite[cell 7]{Notebook}: C, the contradiction matrix, as pairs (i, j), i < j. \\
\texttt{x} & \cite[§5.3, Eq. (4)]{Paper} \cite[cell 11]{Notebook}: x[i, k] = 1 iff exam i is in batch k + 1 \\
\texttt{constraint1} & \cite[§4, known penalties, row 5]{Glover} \cite[§4.2, Eq. (30)]{Lucas}: each exam in at most one batch, the paper's exactly-one \cite[§5.3, Eqs. (5)--(7)]{Paper} relaxed (ADR-0009). With one batch there is no pair to forbid. \\
\texttt{constraint2} & \cite[§5.3, Eqs. (8)--(9)]{Paper} \cite[cell 12]{Notebook}: conflicting exams never share a batch, the penalty x[i, k] x[j, k] of \cite[§4, known penalties, row 1]{Glover}. \\
\texttt{profit} & \cite[§4.2]{Glover} \cite[§4.2, Eq. (31)]{Lucas}: scheduling exam i earns its weight; an unscheduled exam costs it \\
\texttt{A, H, model} & \cite[§5.3, Eq. (10)]{Paper} \cite[cell 13]{Notebook}: H, with lambda = gamma = A and the profit added; A is fed from the Solver Configuration \\
\midrule
\multicolumn{2}{@{}l}{\textbf{Python}, \texttt{examschedule.py}: the runtime the generated scripts call} \\
\texttt{schedule\_of} & \cite[§5.5]{Paper}: the batch of each exam read off X, x[i, k] = 1 iff exam i is in batch k + 1. \\
\texttt{penalty} & A: the largest weight plus one unless set. An explicit A must exceed every weight, or breaking a constraint can cost less than leaving an exam unscheduled (ADR-0009). \cite[§4.2]{Lucas}: "Choosing B < A ensures that it is never favorable to violate the constraint", per weight here. \cite[§4, Comment on the Scalar Penalty P]{Glover}: P large enough that a hard constraint is never broken. \\
\texttt{anneal} & \cite[cell 15]{Notebook}: simulated annealing (neal) on the BQM, then pyqubo decodes every read. \\
\texttt{qaoa} & \cite[Eqs. (1)--(8)]{Farhi}: QAOA on the Ising form of the compiled model, simulated on Aer, with C = H and the expectation minimised, where Farhi maximises it. \\
\texttt{qaoa} & \cite[§2.3, Eq. (12)]{Lucas}: pyqubo's spins, x = (s + 1)/2, so s = 2x - 1. \\
\texttt{qaoa} & \cite[Eq. (5)]{Farhi}: |s>, the uniform superposition, top eigenstate of the mixer B \\
\texttt{qaoa} & \cite[Eq. (6)]{Farhi}: p alternating cost and mixer layers. \\
\texttt{qaoa} & \cite[Eq. (2)]{Farhi}: cost layer U(C, gamma) = exp(-i gamma C), C = H/scale. \cite[§1.1, Eq. (4)]{Lucas}: on a qubit Z = 1 - 2x = -s, so a field h on s acts as -h on Z. \\
\texttt{qaoa} & \cite[Eqs. (3)--(4)]{Farhi}: mixer layer U(B, beta) = exp(-i beta B), B = sum\_j X\_j. \\
\texttt{expected\_energy} & \cite[Eq. (7)]{Farhi}: F\_p, the expectation of the cost in the state, estimated from the shots. \\
\texttt{qaoa} & \cite[§VI]{Farhi}: a linear ramp (gamma grows, |beta| shrinks) on the Trotterised adiabatic path from |s>, the top of B, to the top of -C, the minimum of C; with the cost layer exp(-i gamma C), that path has gamma and beta of opposite signs. \\
\texttt{qaoa} & \cite[Eq. (8)]{Farhi}: the classical loop optimises the 2p angles. \\
\bottomrule
\end{tabularx}
\end{table}
